% Part: incompleteness
% Chapter: introduction
% Section: decidability

\documentclass[../../../include/open-logic-section]{subfiles}

\begin{document}

\olfileid{inc}{int}{dec}

\olsection{अनिर्णेयता आणि अपूर्णता}

ग्योडेलच्या अपूर्णता प्रमेयांच्या सिद्धतेसाठी विन्यासमीमांसेचे
अंकगणितीकरण आवश्यक आहे. पण ते न करताही, एखादी उपपत्ती सर्व
!!{decidable} संबंध !!{represents} करते या गृहीतकावरून काही
चांगले निष्कर्ष मिळतात. ही सिद्धता थांबण्याच्या समस्येची अनिर्णेयता
सिद्ध करणाऱ्या विकर्ण युक्तिवादासारखी आहे.

\begin{thm}
$\Gamma$ ही प्रत्येक !!{decidable} संबंध !!{represents}
करणारी सुसंगत उपपत्ती असेल, तर $\Gamma$ !!{decidable} नाही.
\end{thm}

\begin{proof}
$\Gamma$ !!{decidable} आहे असे समजू. $\Gamma$ प्रत्येक
!!{decidable} संबंध !!{represents} करत असेल तर विसंगत असलीच
पाहिजे, हे दाखवू.

!!^{decidable} गुणधर्म (एकस्थानी संबंध) एक मुक्त चल असलेल्या
!!{formula}s ने निरूपित होतात. अशा सर्व !!{formula}s चे
$!A_0(x)$, $!A_1(x)$, \dots, हे संगणनक्षम प्रगणन असू द्या.
आता पुढील संच $D \subseteq \Nat$ पाहा:
\[
D = \Setabs{n}{\Gamma \Proves \lnot !A_n(\num{n})}
\]
$D$ हा !!{decidable} संच आहे. कारण प्रथम $!A_n(x)$ संगणित
करून, त्यावरून $\lnot !A_n(\num{n})$ मिळवून $n \in D$
आहे का हे तपासता येते. $!A_n(x)$ मधील $x$ च्या प्रत्येक मुक्त
आस्थितीऐवजी पद $\num{n}$ ठेवणे आणि $!A_n(\num{n})$ च्या
आधी $\lnot$ लावणे ही यांत्रिक प्रक्रिया आहे. गृहीतकानुसार
$\Gamma$ !!{decidable} असल्यामुळे
$\lnot !A_n(\num{n}) \in \Gamma$ आहे का हे तपासता येते.
असल्यास $n \in D$, नसल्यास $n \notin D$. त्यामुळे $D$
देखील !!{decidable} आहे.

$\Gamma$ सर्व !!{decidable} गुणधर्म !!{represents} करत
असल्याने ती~$D$ देखील !!{represents} करते. $\Gamma$ मध्ये
$D$ !!{represents}s करणारी सर्व !!{formula}s ही $!A_0(x)$,
$!A_1(x)$, \dots, या यादीत आहेत. म्हणून $!A_d(x)$ हे
$\Gamma$ मध्ये $D$ !!{represents} करेल असा क्रमांक $d$
घ्या. $d
\notin D$ असेल, तर $!A_d(x)$ हे~$D$ !!{represents}
करत असल्यामुळे $\Gamma \Proves
\lnot !A_d(\num{d})$. पण मग $d$ हा~$D$ च्या व्याख्येतील
अट पूर्ण करतो, म्हणून $d \in D$. हे $d \notin D$ शी
विसंगत आहे. म्हणून अप्रत्यक्ष सिद्धतेने $d \in D$.

$d \in D$ असल्यामुळे~$D$ च्या व्याख्येनुसार $\Gamma \Proves \lnot
!A_d(\num{d})$. दुसरीकडे $!A_d(x)$ हे $\Gamma$ मध्ये~$D$
!!{represents} करत असल्यामुळे $\Gamma \Proves !A_d(\num{d})$.
म्हणून $\Gamma$ विसंगत आहे.
\end{proof}

\begin{explain}
वरील प्रमेय दाखवते की सर्व !!{decidable} संबंध !!{represents}
करणारी कोणतीही सुसंगत उपपत्ती !!{decidable} असू शकत नाही.
$\Th{Q}$ सर्व !!{decidable} संबंध !!{represents}s करते हे
आपण दाखवू. म्हणजे $\Th{Q}$ समाविष्ट असलेल्या $\Th{PA}$
आणि $\Th{TA}$ यांसारख्या सर्व उपपत्तीदेखील तसे करतात; त्यांपैकी
सुसंगत उपपत्ती !!{decidable} नसतात. (येथे नाव घेतलेल्या उपपत्ती
प्रमाण प्रतिमानात सत्य असल्यामुळे सुसंगत आहेत.)

या निष्कर्षाचा वापर करून पहिल्या अपूर्णता प्रमेयाची दुर्बल आवृत्तीही
मिळते. !!{axiomatizable} आणि !!{complete} असलेली कोणतीही
उपपत्ती !!{decidable} असते. त्यामुळे सुसंगत, !!{axiomatizable}
आणि सर्व !!{decidable} गुणधर्म !!{represents}s करणारी उपपत्ती
!!{complete} असू शकत नाही.
\end{explain}

\begin{thm}
$\Gamma$ !!{axiomatizable} आणि !!{complete} असेल, तर ती
!!{decidable} असते.
\end{thm}

\begin{proof}
कोणतीही विसंगत उपपत्ती !!{decidable} असते, कारण विसंगत
उपपत्त्यांत सर्व !!{sentence}s असतात. म्हणून ``$!A \in
\Gamma$ आहे का?'' या प्रश्नाचे उत्तर नेहमी ``होय'' असते,
म्हणजे तो निर्णय करता येतो.

आता $\Gamma$ सुसंगत, !!{axiomatizable} आणि !!{complete} आहे
असे समजू. $\Gamma$ !!{axiomatizable} असल्याने ती
!!{computably enumerable} आहे. कारण~$\Gamma$ च्या
स्वयंसिद्धकांपासूनच्या सर्व योग्य !!{derivation}s चे संगणनक्षम
फलनाने प्रगणन करता येते. योग्य !!{derivation} पासून ती ज्या
!!{sentence} ची !!{derive}s करते ते संगणित करता येते. त्यामुळे
दोन्ही मिळून~$\Gamma$ च्या सर्व प्रमेयांचे प्रगणन करणारे
संगणनक्षम फलन मिळते. $\Gamma$ सुसंगत आणि !!{complete} असल्यामुळे
एखादे !!{sentence} हे~$\Gamma$ चे प्रमेय असते तेव्हाच
$\lnot !A$ प्रमेय नसते. म्हणून $!A \in
\Gamma$ आहे का हे पुढीलप्रमाणे ठरवता येते. $\Gamma$ च्या सर्व
प्रमेयांचे प्रगणन करा. या यादीत $!A$ आल्यास $\Gamma \Proves !A$
हे कळते. यादीत $\lnot
!A$ आल्यास $\Gamma \Proves/ !A$ हे कळते. $\Gamma$
!!{complete} असल्यामुळे यांपैकी एक प्रकरण अखेरीस घडतेच;
म्हणून प्रक्रिया अखेरीस उत्तर देते.
\end{proof}

\begin{cor}
\ollabel{cor:incompleteness}
$\Gamma$ सुसंगत, !!{axiomatizable} आणि प्रत्येक !!{decidable}
गुणधर्म !!{represents} करणारी असेल, तर ती !!{complete} नाही.
\end{cor}

\begin{proof}
$\Gamma$ !!{complete} असेल, तर मागील प्रमेयामुळे ती
!!{decidable} असेल (कारण ती !!{axiomatizable} आणि सुसंगत आहे).
पण $\Gamma$ प्रत्येक !!{decidable} गुणधर्म !!{represents}
करत असल्यामुळे पहिल्या प्रमेयाने ती !!{decidable} नाही.
\end{proof}

\begin{prob}
$\Th{TA} = \Setabs{!A}{\Sat{N}{!A}}$ ही !!{axiomatizable}
नाही, हे दाखवा. $\Th{TA}$ सर्व निर्णेय गुणधर्मांचे निरूपण करते
असे गृहीत धरू शकता.
\end{prob}

उदाहरणार्थ, $\Th{Q}$ सर्व !!{decidable} गुणधर्म !!{represents}
करते हे आपण प्रस्थापित केल्यावर उपप्रमेयावरून $\Th{Q}$ अपूर्ण
असलीच पाहिजे असे कळेल. पण त्याची सिद्धता स्वतंत्र !!{sentence}
चे उदाहरण देत नाही; असे !!a{sentence} अस्तित्वात असलेच पाहिजे
एवढेच दाखवते. त्यासाठी विन्यासमीमांसेचे अंकगणितीकरण करून
ग्योडेलच्या मूळ सिद्धतेची कल्पना वापरावी लागेल. आणि अर्थातच,
$\Th{Q}$ सर्व !!{decidable} गुणधर्म !!{represents}s करते हा
पहिला दावाही अजून सिद्ध करायचा आहे.

प्रत्येक \emph{रंजक} उपपत्ती अपूर्ण किंवा अनिर्णेय असेलच असे
नाही. अनेक उपपत्ती महत्त्वाची गणितीय तथ्ये व्यक्त करण्याइतक्या
प्रबळ असूनही ग्योडेलच्या निष्कर्षांच्या अटी पूर्ण करत नाहीत.
उदाहरणार्थ, $\Th{Pres} = \Setabs{!A \in \Lang{L_{A^+}}}{\Sat{N}{!A}}$
ही प्रमाण प्रतिमानात सत्य असलेल्या, पण~$\times$ नसलेल्या
अंकगणिताच्या भाषेतील !!{sentence}s ची उपपत्ती संपूर्ण आणि
निर्णेय दोन्ही आहे. या उपपत्तीला प्रेस्बर्गर अंकगणित म्हणतात;
केवळ $\Obj{0}$, $\prime$ आणि~$+$ यांच्या साहाय्याने मांडता
येणारी नैसर्गिक संख्यांविषयीची सर्व सत्ये ती सिद्ध करते.

\end{document}
